\documentclass[11pt]{article}
\usepackage{amsmath,amssymb,amsthm}
\usepackage{algorithm2e}
\usepackage{enumitem}
\usepackage{fullpage}
\usepackage{hyperref}
\newtheorem{theorem}{Theorem}
\newtheorem{lemma}{Lemma}
\newcommand{\EE}{\mathbb{E}}
\newcommand{\RR}{\mathbb{R}}
\newcommand{\norm}[1]{\left\lVert#1\right\rVert}
\newcommand{\ip}[2]{\left\langle #1,#2\right\rangle}
\newcommand{\e}{\mathbf e}
\begin{document}

\section*{Randomized Coordinate Descent with Probability–Weighted Updates (RCD‑PW)}

\section*{Mathematical Formulation}
Let $f:\RR^{d}\rightarrow\RR$ be a convex differentiable function.  
Denote by $e_i\in\RR^{d}$ the $i$‑th canonical basis vector.  
At iteration $k$ a coordinate $i_k\in\{1,\dots ,d\}$ is drawn independently
according to a fixed probability distribution 
\[
\Pr(i_k=i)=p_i,\qquad p_i>0,\ \sum_{i=1}^{d}p_i=1 .
\]

Define the \emph{probability‑weighted stochastic gradient}
\[
g_k\;:=\;\frac{1}{p_{i_k}}\;\partial_i f(w_k)\;e_{i_k},
\qquad\text{where }\partial_i f(w_k)=\big\langle\nabla f(w_k),e_i\big\rangle .
\tag{1}
\]

The update rule is
\[
w_{k+1}=w_k-\eta\,g_k,
\qquad\eta>0\ \text{(stepsize)} .
\tag{2}
\]

Because of the factor $1/p_{i_k}$ we have
\[
\EE\!\left[g_k\mid w_k\right]=\nabla f(w_k),
\tag{3}
\]
hence $g_k$ is an unbiased estimator of the full gradient.

\section*{Pseudocode}
\begin{algorithm}[H]
\SetAlgoLined
\KwIn{Initial point $w_0\in\RR^{d}$, stepsize $\eta>0$, probabilities $\{p_i\}_{i=1}^{d}$}
\For{$k=0,1,\dots,T-1$}{
   Sample $i_k\sim\{1,\dots ,d\}$ with $\Pr(i_k=i)=p_i$\;
   Compute partial derivative $g^{(i_k)}_k:=\partial_{i_k}f(w_k)$\;
   Form the weighted coordinate gradient $g_k:=\dfrac{g^{(i_k)}_k}{p_{i_k}}\,e_{i_k}$\;
   Update $w_{k+1}\gets w_k-\eta\,g_k$\;
}
\KwOut{Final iterate $w_T$ (or averaged iterate $\bar w_T$).}
\caption{Randomized Coordinate Descent with Probability‑Weighted Updates}
\end{algorithm}

\section*{Dry Run Example}
Consider the one‑dimensional quadratic
\[
f(w)=(w-3)^{2},\qquad \nabla f(w)=2(w-3).
\]
We have $d=1$, $p_{1}=1$, and $L_{1}=2$ (the Lipschitz constant of $\partial_{1}f$).  
Take $w_{0}=0$, stepsize $\eta=0.1$.

\begin{center}
\begin{tabular}{c|c|c|c}
Iteration $k$ & $g_k$ (by (1)) & $w_{k+1}=w_k-\eta g_k$ & $f(w_{k+1})$\\\hline
0 & $\dfrac{2(w_0-3)}{1}=2(-3)=-6$ & $0-0.1(-6)=0.6$ & $(0.6-3)^2=5.76$\\
1 & $2(w_1-3)=2(0.6-3)=-4.8$ & $0.6-0.1(-4.8)=1.08$ & $(1.08-3)^2=3.6864$\\
2 & $2(w_2-3)=2(1.08-3)=-3.84$ & $1.08-0.1(-3.84)=1.464$ & $(1.464-3)^2=2.3665$\\
\end{tabular}
\end{center}
The objective value decreases at each step, illustrating the algorithm.

\newpage

\section*{Assumptions}
\begin{enumerate}[label=\textbf{A\arabic*}, leftmargin=*]
\item \textbf{Convexity.} $f$ is convex:
\[
f(y)\ge f(x)+\langle\nabla f(x),y-x\rangle,\qquad\forall x,y\in\RR^{d}.
\tag{A1}
\]

\item \textbf{Coordinate‑wise $L_i$‑smoothness.} For each $i\in\{1,\dots ,d\}$ there exists $L_i>0$ such that
\[
\big|\partial_i f(x+\alpha e_i)-\partial_i f(x)\big|\le L_i|\alpha|,
\qquad\forall x\in\RR^{d},\ \forall\alpha\in\RR .
\tag{A2}
\]
Equivalently,
\[
f(x+\alpha e_i)\le f(x)+\partial_i f(x)\alpha+\frac{L_i}{2}\alpha^{2}.
\tag{A2'}
\]

\item \textbf{Existence of a minimizer.} The set of minimizers
\[
\mathcal W^{*}:=\arg\min_{w\in\RR^{d}}f(w)
\]
is non‑empty; denote any $w^{*}\in\mathcal W^{*}$.

\item \textbf{Stepsize restriction.} The stepsize satisfies
\[
0<\eta\le\frac{1}{\displaystyle\max_{i=1,\dots ,d}\frac{L_i}{p_i}} .
\tag{A3}
\]
\end{enumerate}

\section*{Main Convergence Theorem}
\begin{theorem}[Expected suboptimality of RCD‑PW]\label{thm:main}
Under Assumptions \textbf{A1}–\textbf{A3}, the iterates generated by
\eqref{2} satisfy for every $T\ge1$
\[
\EE\!\big[\,f(w_T)-f(w^{*})\,\big]\;\le\;
\frac{\norm{w_0-w^{*}}^{2}}{2\eta T}.
\tag{4}
\]
Consequently, the averaged iterate $\bar w_T:=\frac1T\sum_{k=0}^{T-1}w_k$ enjoys the same bound because of convexity.
\end{theorem}

\section*{Theoretical Properties}
\begin{enumerate}[label=\textbf{P\arabic*}, leftmargin=*]
\item \textbf{Linear speed‑up with respect to the probability weights.}
If a coordinate $i$ is sampled more often (larger $p_i$), the factor $L_i/p_i$ in (A3) becomes smaller, allowing a larger stepsize for that direction.

\item \textbf{Uniform sampling as a special case.}
When $p_i=1/d$ for all $i$, condition (A3) reduces to $\eta\le d/\max_i L_i$, reproducing the classical bound for uniform randomized coordinate descent.

\item \textbf{Stability.}
The bound (4) is independent of the variance of the stochastic gradient because $g_k$ is an unbiased estimator with a deterministic second‑moment bound derived from (A2).

\item \textbf{Comparison to full‑gradient GD.}
If $p_i\equiv1$ (i.e., always pick the same coordinate) the method reduces to a one‑dimensional gradient descent on that coordinate, and (4) coincides with the standard GD bound with stepsize $\eta\le 1/L_i$.
\end{enumerate}

\section*{Discussion}
The theorem shows an $\mathcal O(1/T)$ convergence rate in expectation, identical to that of full‑gradient gradient descent for convex smooth functions. The probability‑weighted scaling renders the stochastic direction an unbiased estimator of the full gradient, which is the key element enabling the simple $1/T$ rate without extra variance terms. Choosing $p_i$ proportional to $L_i$ (i.e., $p_i\propto L_i$) balances the term $L_i/p_i$ across coordinates and maximises the admissible stepsize, yielding the fastest practical convergence.

\newpage

\section*{Preliminaries (Definitions and Lemmas)}
\begin{lemma}[Coordinate smoothness inequality]\label{lem:smooth}
For any $x\in\RR^{d}$, any coordinate $i$, and any scalar $\alpha$,
\[
f(x+\alpha e_i)\le f(x)+\partial_i f(x)\,\alpha+\frac{L_i}{2}\alpha^{2}.
\]
\end{lemma}
\begin{proof}
This is exactly Assumption (A2'). \qedhere
\end{proof}

\begin{lemma}[Unbiasedness of the weighted gradient]\label{lem:unbiased}
For the definition \eqref{1},
\[
\EE\!\left[g_k\mid w_k\right]=\nabla f(w_k).
\]
\end{lemma}
\begin{proof}
\[
\EE\!\left[g_k\mid w_k\right]
=\sum_{i=1}^{d}p_i\frac{1}{p_i}\partial_i f(w_k)e_i
=\sum_{i=1}^{d}\partial_i f(w_k)e_i
=\nabla f(w_k).
\qquad\qedhere
\]
\end{proof}

\section*{Main Proof (Step‑by‑Step Derivation)}
We prove Theorem \ref{thm:main}. Throughout we fix $w^{*}\in\mathcal W^{*}$.

\noindent\textbf{Step 1 (Descent Lemma for a single iteration).}\\
Using Lemma \ref{lem:smooth} with $\alpha=-\eta\,\frac{1}{p_{i_k}}\partial_{i_k}f(w_k)$ we obtain
\begin{align}
f(w_{k+1})
&=f\!\big(w_k-\eta g_k\big) \nonumber\\
&\le f(w_k)-\eta\;\partial_{i_k}f(w_k)\,\frac{1}{p_{i_k}}\partial_{i_k}f(w_k)
   +\frac{L_{i_k}}{2}\Big(\eta\frac{1}{p_{i_k}}\partial_{i_k}f(w_k)\Big)^{2} \nonumber\\
&=f(w_k)-\eta\,\frac{(\partial_{i_k}f(w_k))^{2}}{p_{i_k}}
   +\frac{\eta^{2}L_{i_k}}{2}\frac{(\partial_{i_k}f(w_k))^{2}}{p_{i_k}^{2}} .
\tag{5}
\end{align}
\hfill[Step 1]

\noindent\textbf{Step 2 (Take conditional expectation).}\\
Condition on $w_k$ and use $\Pr(i_k=i)=p_i$:
\begin{align}
\EE\!\big[f(w_{k+1})\mid w_k\big]
&\le f(w_k)-\eta\sum_{i=1}^{d}p_i\frac{(\partial_i f(w_k))^{2}}{p_i}
   +\frac{\eta^{2}}{2}\sum_{i=1}^{d}p_i\frac{L_i}{p_i^{2}}(\partial_i f(w_k))^{2} \nonumber\\
&= f(w_k)-\eta\|\nabla f(w_k)\|^{2}
   +\frac{\eta^{2}}{2}\sum_{i=1}^{d}\frac{L_i}{p_i}(\partial_i f(w_k))^{2}.
\tag{6}
\end{align}
\hfill[Step 2]

\noindent\textbf{Step 3 (Relate gradient norm to distance from optimum).}\\
Convexity (A1) with $y=w^{*}$ gives
\[
f(w_k)-f(w^{*})\le\langle\nabla f(w_k),w_k-w^{*}\rangle .
\tag{7}
\]
Apply Cauchy–Schwarz:
\[
f(w_k)-f(w^{*})\le \|\nabla f(w_k)\|\,\|w_k-w^{*}\|.
\tag{8}
\]

\noindent\textbf{Step 4 (Quadratic expansion of the distance).}\\
From the update rule \eqref{2},
\begin{align}
\|w_{k+1}-w^{*}\|^{2}
&= \|w_k-w^{*}-\eta g_k\|^{2} \nonumber\\
&= \|w_k-w^{*}\|^{2}
   -2\eta\langle g_k,w_k-w^{*}\rangle
   +\eta^{2}\|g_k\|^{2}.
\tag{9}
\end{align}
\hfill[Step 4]

\noindent\textbf{Step 5 (Take conditional expectation of (9)).}\\
Using Lemma \ref{lem:unbiased} and $\EE[\|g_k\|^{2}\mid w_k]
      =\sum_{i=1}^{d}\frac{1}{p_i^{2}}(\partial_i f(w_k))^{2}p_i
      =\sum_{i=1}^{d}\frac{(\partial_i f(w_k))^{2}}{p_i}$, we obtain
\begin{align}
\EE\!\big[\|w_{k+1}-w^{*}\|^{2}\mid w_k\big]
&= \|w_k-w^{*}\|^{2}
   -2\eta\langle\nabla f(w_k),w_k-w^{*}\rangle
   +\eta^{2}\sum_{i=1}^{d}\frac{(\partial_i f(w_k))^{2}}{p_i}.
\tag{10}
\end{align}
\hfill[Step 5]

\noindent\textbf{Step 6 (Combine (6) and (10) to eliminate the gradient norm).}\\
From (8) we have $\langle\nabla f(w_k),w_k-w^{*}\rangle\ge f(w_k)-f(w^{*})$.
Insert this lower bound into (10):
\begin{align}
\EE\!\big[\|w_{k+1}-w^{*}\|^{2}\mid w_k\big]
&\le \|w_k-w^{*}\|^{2}
   -2\eta\big(f(w_k)-f(w^{*})\big)
   +\eta^{2}\sum_{i=1}^{d}\frac{(\partial_i f(w_k))^{2}}{p_i}.
\tag{11}
\end{align}
\hfill[Step 6]

\noindent\textbf{Step 7 (Upper bound the last term via (6)).}\\
Rearrange (6) as
\[
\eta\|\nabla f(w_k)\|^{2}
\le f(w_k)-\EE\!\big[f(w_{k+1})\mid w_k\big]
   +\frac{\eta^{2}}{2}\sum_{i=1}^{d}\frac{L_i}{p_i}(\partial_i f(w_k))^{2}.
\]
Multiply by $2\eta$ and add $\eta^{2}\sum_{i}\frac{(\partial_i f(w_k))^{2}}{p_i}$ to both sides:
\begin{align}
2\eta^{2}\|\nabla f(w_k)\|^{2}
&\le 2\eta\big(f(w_k)-\EE[f(w_{k+1})\mid w_k]\big)
   +\eta^{3}\sum_{i=1}^{d}\frac{L_i}{p_i}(\partial_i f(w_k))^{2}
   +\eta^{2}\sum_{i=1}^{d}\frac{(\partial_i f(w_k))^{2}}{p_i}.
\tag{12}
\end{align}
Since $\eta\le 1/\max_i(L_i/p_i)$ (Assumption A3), we have
\[
\eta^{3}\frac{L_i}{p_i}\le\eta^{2}\frac{1}{p_i},
\]
hence the two sums on the right–hand side of (12) can be bounded by
$2\eta^{2}\sum_{i}\frac{(\partial_i f(w_k))^{2}}{p_i}$.
Thus
\[
\eta^{2}\sum_{i=1}^{d}\frac{(\partial_i f(w_k))^{2}}{p_i}
\le \eta\big(f(w_k)-\EE[f(w_{k+1})\mid w_k]\big).
\tag{13}
\]
\hfill[Step 7]

\noindent\textbf{Step 8 (Plug (13) into (11)).}\\
Using (13) in (11) gives
\begin{align}
\EE\!\big[\|w_{k+1}-w^{*}\|^{2}\mid w_k\big]
&\le \|w_k-w^{*}\|^{2}
   -2\eta\big(f(w_k)-f(w^{*})\big)
   +\eta\big(f(w_k)-\EE[f(w_{k+1})\mid w_k]\big) \nonumber\\
&= \|w_k-w^{*}\|^{2}
   -\eta\big(f(w_k)-f(w^{*})\big)
   -\eta\big(\EE[f(w_{k+1})\mid w_k]-f(w^{*})\big).
\tag{14}
\end{align}
\hfill[Step 8]

\noindent\textbf{Step 9 (Take total expectation).}\\
Denote $\Delta_k:=\EE\big[f(w_k)-f(w^{*})\big]$.
Taking expectation of (14) over all randomness up to iteration $k$ yields
\[
\EE\!\big[\|w_{k+1}-w^{*}\|^{2}\big]
\le \EE\!\big[\|w_k-w^{*}\|^{2}\big]
      -\eta\Delta_k-\eta\Delta_{k+1}.
\tag{15}
\]
Rearranging,
\[
\eta\Delta_{k+1}\le \EE\!\big[\|w_k-w^{*}\|^{2}\big]
                     -\EE\!\big[\|w_{k+1}-w^{*}\|^{2}\big]
                     -\eta\Delta_k .
\tag{16}
\]
\hfill[Step 9]

\noindent\textbf{Step 10 (Telescoping sum).}\\
Sum \eqref{16} for $k=0,\dots,T-1$:
\[
\eta\sum_{k=0}^{T-1}\Delta_{k+1}
\le \|w_0-w^{*}\|^{2}
    -\EE\!\big[\|w_T-w^{*}\|^{2}\big]
    -\eta\sum_{k=0}^{T-1}\Delta_k .
\]
Move the term $-\eta\sum_{k=0}^{T-1}\Delta_k$ to the left:
\[
\eta\big(\Delta_T+\Delta_{T-1}+\dots+\Delta_1\big)
      +\eta\big(\Delta_{T-1}+\dots+\Delta_0\big)
\le \|w_0-w^{*}\|^{2}.
\]
Hence $2\eta\sum_{k=1}^{T}\Delta_k\le \|w_0-w^{*}\|^{2}$, i.e.
\[
\frac{1}{T}\sum_{k=1}^{T}\Delta_k\le\frac{\|w_0-w^{*}\|^{2}}{2\eta T}.
\tag{17}
\]
Since $f$ is convex, the average iterate $\bar w_T$ satisfies
$f(\bar w_T)-f(w^{*})\le\frac1T\sum_{k=1}^{T}\Delta_k$,
and therefore
\[
\EE\!\big[f(\bar w_T)-f(w^{*})\big]\le\frac{\|w_0-w^{*}\|^{2}}{2\eta T}.
\]
Because $\EE[f(w_T)]\le\EE[f(\bar w_T)]$, the same bound holds for the last iterate, completing the proof of \eqref{4}.
\hfill[Step 10] \qed

\section*{Rate Derivation (With Constants)}
The bound \eqref{4} can be rewritten as
\[
\EE\!\big[f(w_T)-f(w^{*})\big]\le
\underbrace{\frac{\|w_0-w^{*}\|^{2}}{2}}_{C_0}\,
\underbrace{\frac{1}{\eta}}_{C_1}\,\frac{1}{T}.
\]
Choosing the maximal admissible stepsize from (A3),
\[
\eta^{\star}= \frac{1}{\displaystyle\max_{i}\frac{L_i}{p_i}},
\]
gives the explicit constant
\[
C_1 = \max_{i}\frac{L_i}{p_i}.
\]
Thus
\[
\EE\!\big[f(w_T)-f(w^{*})\big]\le
\frac{\max_{i}\frac{L_i}{p_i}\,\|w_0-w^{*}\|^{2}}{2\,T}.
\]

\section*{Special Cases}
\begin{enumerate}[label=\arabic*.]
\item \textbf{Uniform sampling.} $p_i=1/d$. Then $\max_i\frac{L_i}{p_i}=d\max_i L_i$ and the stepsize condition becomes $\eta\le 1/(d\max_i L_i)$. The bound reduces to the classical $\mathcal O(dL_{\max}\|w_0-w^{*}\|^{2}/T)$ rate for uniform randomized coordinate descent.

\item \textbf{Probability proportional to smoothness.} Set $p_i = L_i/\sum_{j}L_j$. Then $\frac{L_i}{p_i}= \sum_{j}L_j$ for all $i$, so the stepsize bound simplifies to $\eta\le 1/\sum_{j}L_j$, and the convergence constant becomes $\big(\sum_{j}L_j\big)\|w_0-w^{*}\|^{2}/(2T)$, which is often much tighter than the uniform case when the $L_i$ are heterogeneous.

\item \textbf{Single coordinate descent.} If $p_{i^{\star}}=1$ for some $i^{\star}$ and $p_i=0$ otherwise, the algorithm coincides with ordinary gradient descent on the $i^{\star}$‑th coordinate, and the bound recovers the one‑dimensional smooth convex rate $\|w_0-w^{*}\|^{2}/(2\eta T)$ with $\eta\le 1/L_{i^{\star}}$.
\end{enumerate}

\end{document}